\documentclass[english, parskip=full]{scrartcl}
\usepackage[utf8]{inputenc}  % Kodierung der Datei
\usepackage[ngerman]{babel}  % Sprache auf Deutsch 
\usepackage{color}
\usepackage{graphicx, gensymb}
\usepackage{amsfonts, cancel, amssymb}
\usepackage{amsmath, amsthm, array, esint}
\usepackage{booktabs, multirow}  % schönere Tabellen
\usepackage{caption}  % Anpassung der Captions
\usepackage{scrlayer-scrpage}    % Manipulation des Seitenstils
%\pagestyle{scrheadings}  % Stil für die Seite setzen
%\automark{section}  % Abschnittsnamen als Seitenbeschriftung 
%\setheadsepline{.5pt}    % Kopzeile durch Linie abtrennen
\usepackage[hidelinks]{hyperref}  % Links und weitere PDF-Features
\usepackage{typearea, tikz, pgffor, verbatim}
\usetikzlibrary{calc, quotes}
\usepackage[cache=false, outputdir=build]{minted}
\usepackage{placeins} % provides \FloatBarrier

\definecolor{bg}{rgb}{0.9,0.9,0.9}
\newcommand\numberthis{\addtocounter{equation}{1}\tag{\theequation}}
\usepackage{siunitx}
\newcommand{\bra}{}
\newcommand{\kom}{}
\newcommand{\brak}{}
\newcommand{\braket}{}
\newcommand{\intx}{}
\newcommand{\intr}{}
\newcommand{\kla}{}
\newcommand{\klam}{}
\newcommand{\klammer}{}
\newcommand{\oper}{}
\newcommand{\tecxt}{}
\newcommand{\Bra}{}
\newcommand{\Ket}{}
\renewcommand{\i}{\text{i}}
\newcommand{\e}{\text{e}}
\newcommand*{\Fourier}{\textsc{Fourier}\ }
\newcommand*{\F}{\mathcal{F}}
\newcommand*{\sinc}{\text{sinc\,}}
\newcommand{\conv}{\mathop{\scalebox{3.5}{\raisebox{-0.38ex}{\makebox[0.5\width][c]{$\ast$}}}}\limits}

\newcommand*{\titel}{06 Verschiedene Runge-Kutta Verfahren}
\newcommand*{\autor}{Joachim Schwardt und Julian Fleck}

\hypersetup{pdfauthor={\autor}, pdftitle={\titel}}

%\titlehead{titlehead}
\subject{Numerik II - Aufgaben}
\title{\titel}
%\author{\autor}
\author{\begin{tabular}{l|l}
Joachim Schwardt & 4768711 \\ 
Julian Fleck & 4759587\\ 
\end{tabular}}
\date{\today}

\begin{document}

%\begin{titlepage}
\maketitle  % Titel setzen
%\tableofcontents  % Inhaltsverzeichnis setzen
%\end{titlepage}

\section*{Aufgabe 1}
Gegeben ist ein explizites \textsc{Runge-Kutta}-Verfahren mit dem \textsc{Butcher}-Schema \autoref{table:butcher scheme order 4}.
\begin{table}[!ht]
\centering
\begin{tabular}{c|cccc}
0 & &&&\\[2pt]
$\frac{1}{3}$ & $\frac{1}{3}$ &&&\\[2pt]
$\frac{2}{3}$ & $-\frac{1}{3}$ & 1 &&\\[2pt]
1 & 1 & -1 & 1 &\\[2pt] \hline\\[-10pt]
& $\frac{1}{8}$ &$\frac{3}{8}$ &$\frac{3}{8}$ &$\frac{1}{8}$ \\[2pt]
\end{tabular}
\caption{Butcher-Schema für ein RK-Verfahren.}
\label{table:butcher scheme order 4}
\end{table}\\
Ausführlich entspricht dies
\begin{align}
k_1 &= f\left( t, x \right),\\
k_2 &= f\left( t + \frac{\tau}{3}, x + \tau \frac{k_1}{3} \right),\\
k_3 &= f\left( t + \frac{2\tau}{3}, x + \tau \left[ k_2 - \frac{k_1}{3} \right] \right),\\
k_4 &= f\left( t + \tau, x + \tau \left[ k_1 - k_2 + k_3 \right] \right),\\
\Psi^{t+\tau,t}x &= x + \frac{\tau}{8} \klam\left[ k_1 + 3k_2 + 3k_3 + k_4 \right].
\end{align}
Beziehungsweise noch ausführlicher
\begin{align*}
\Psi^{t+\tau,t}x &= x + \frac{\tau}{8} \Bigg[ f\left( t, x \right) + 3f\left( t + \frac{\tau}{3}, x + \tau \frac{f\left( t, x \right)}{3} \right) + \\
&\hspace{1.7cm} + 3f\left( t + \frac{2\tau}{3}, x + \tau \left[ f\left( t + \frac{\tau}{3}, x + \tau \frac{f\left( t, x \right)}{3} \right) - \frac{f\left( t, x \right)}{3} \right] \right) + \\
&\hspace{1.7cm} + f\Bigg( t + \tau, \\
&\hspace{2.7cm} x + \tau \bigg[ f\left( t, x \right) - f\left( t + \frac{\tau}{3}, x + \tau \frac{f\left( t, x \right)}{3} \right) + \\
&\hspace{3.7cm}+ f\left( t + \frac{2\tau}{3}, x + \tau \left[ f\left( t + \frac{\tau}{3}, x + \tau \frac{f\left( t, x \right)}{3} \right) - \frac{f\left( t, x \right)}{3} \right] \right) \bigg] \\
&\hspace{2.7cm} \Bigg) \\
&\hspace{1.7cm} \Bigg].
\end{align*}

\section*{Aufgabe 2}
Gegeben sei ein $s$-stufiges \textsc{Runge-Kutta}-Verfahren mit Koeffizienten $A,b,c$.
Für ein Anfangswertproblem
\begin{align}
x'(t) &= f(t,x(t))
\end{align}
mit $x(t_0) = x_0$ lautet das Verfahren ausführlich dann
\begin{align}
k_j &= f\kla\left( t + c_j\tau, x + \tau\sum_{l=1}^{j-1} A_{jl}k_l \right),\\
\Psi^{t+\tau,t}x &= x + \tau \sum_{j=1}^s b_jk_j.
\end{align}
Wir könnten die Differentialgleichung zuvor allerdings noch autonomisieren.
Definiere dazu $s(t) := t$ und schreibe $z := (s,x)^\top \in\mathbb{R}^{d+1}$
\begin{align}
z'(t) &= \begin{pmatrix}
s'(t) \\ x'(t)
\end{pmatrix} = \begin{pmatrix}
1 \\ f(s,x)
\end{pmatrix} =: F(z).
\end{align}
Für diese Differentialgleichung lautet dasselbe Verfahren dann
\begin{align}
\tilde{k}_j &= F\kla\left( z + \tau\sum_{l=1}^{j-1} A_{jl} \tilde{k}_l \right) = \begin{pmatrix}
1 \\ f\kla\left( s + \tau\sum_{l=1}^{j-1} A_{jl} (\tilde{k}_l)_1,\ x + \tau \sum_{l=1}^{j-1} A_{jl} (\tilde{k}_l)_2 \right)
\end{pmatrix}.
\end{align}
Der erste Eintrag im Vektor jeder Stufe ist also immer gleich Eins.
Damit vereinfacht sich das Verfahren zu $(\tilde{k}_j)_1 = 1$ und
\begin{align}
(\tilde{k}_j)_2 &= f\kla\left( s + \tau\sum_{l=1}^{j-1} A_{jl},\ x + \tau \sum_{l=1}^{j-1} A_{jl} (\tilde{k}_l)_2 \right).
\end{align}
Die Evolution ergibt dann
\begin{align}
\tilde{\Psi}^{t+\tau,t}z &= z + \tau \sum_{j=1}^s b_j\tilde{k}_j = \begin{pmatrix}
s + \tau \sum_{j=1}^s b_j \\ x + \tau \sum_{j=1}^s b_j (\tilde{k}_j)_2
\end{pmatrix}.
\end{align}
Zum Vergleich, zuvor hatten wir per Konstruktion $t \rightarrow t+\tau$ und
\begin{align}
\Psi^{t+\tau,t} x &= x + \tau \sum_{j=1}^s b_j k_j.
\end{align}
Damit das Verfahren invariant ist, müssen also die Stufen $k_j$ mit der zweiten Komponente der Stufen $\tilde{k}_j$ übereinstimmen.
Wir fordern also
\begin{align}
f\kla\left( s + \tau\sum_{l=1}^{j-1} A_{jl},\ x + \tau \sum_{l=1}^{j-1} A_{jl} k_l \right) &\overset{!}{=} f \kla\left( t + \tau c_j,\ x + \tau\sum_{l=1}^{j-1}A_{jl} k_l \right),
\end{align}
woraus $\sum_{l=1}^{j-1} A_{jl} = c_j$ folgt.
Außerdem muss allerdings die ``Transformation'' von $t$ mit der von $s$ übereinstimmen.
Da $s \rightarrow s + \tau \sum_{j=1}^s b_j$ gilt, ist das also äquivalent zu $\sum_{j=1}^s b_j = 1$, was wiederum äquivalent dazu ist, dass das Verfahren konsistent ist.

\section*{Aufgabe 3}
Gegeben sei eine stetige Funktion $f : \mathbb{R} \rightarrow \mathbb{R}$ sowie ein Punkt $(t_0,x_0) \in \mathbb{R}^2$.
Wir betrachten das Anfangswertproblem 
\begin{align}
x'(t) &= f(t), \quad x(t_0) = x_0.
\end{align}
In diesem Fall gilt
\begin{align}
\frac{\tecxt\text{d}}{\tecxt\text{d}\tau} \Phi^{t+\tau,t}x = \frac{\tecxt\text{d}}{\tecxt\text{d}\tau} x(t+\tau) = x'(t+\tau) = f(t+\tau).
\end{align}
Insbesondere vereinfacht sich die Taylor-Entwicklung der Evolution damit zu
\begin{align}
\Phi^{t+\tau,t} x &= x + \tau  f + \frac{\tau^2}{2} f' + \frac{\tau^3}{6} f'' + \dots
\end{align}
Es seien noch Gitterpunkte $t_0 < \dots < t_n$ gegeben.
Zur Berechnung von Näherungswerten $x_h(t_j)$ für die exakten Funktionswerte $x(t_j)$ werde ein $s$-stufiges \textsc{Runge-Kutta}-Verfahren mit der Konsistenzordnung $p>0$ verwendet.
\\
Sei $f$ ein Polynom vom Grad höchstens $(p-1)$.
Dann können wir schreiben
\begin{align}
f(t) &= \sum_{\rho=0}^{p-1} w_\rho t^\rho,
\end{align}
wobei $w_\rho$ nicht näher bestimmte Koeffizienten sind.
Die exakte Lösung der Differentialgleichung ist dann einfach
\begin{align}
x(t) &= x(t_0) + \intx\int_{t_0}^{t} f(s)\, \mathrm{d}s = x_0 + \sum_{\rho=0}^{p-1} w_\rho \frac{t^{\rho + 1} - t_0^{\rho + 1}}{\rho + 1} \equiv \Phi^{t,t_0}x(t_0).
\end{align}
Die Evolution können wir also als
\begin{align}
\Phi^{t+\tau,t}x &= x + \sum_{\rho=0}^{p-1} w_\rho \frac{(t+\tau)^{\rho + 1} - t^{\rho + 1}}{\rho + 1}
\end{align}
schreiben.
%Für die Evolution gilt dann exakt
%\begin{align}
%\Phi^{t+\tau,t} x &= x + \tau  f + \frac{\tau^2}{2} f' + \frac{\tau^3}{6} f'' + \dots + \frac{\tau^p}{p!}f^{(p-1)} = \\
%&= x + \sum_{k=1}^p \frac{\tau^k}{k!} \frac{\tecxt\text{d}^{k-1}}{\tecxt\text{d}t^{k-1}} \sum_{\rho=0}^{p-1} w_{\rho} t^{\rho} = \\
%&= x + \sum_{\rho = 0}^{p-1} \sum_{k=1}^p w_\rho \frac{\rho (\rho-1)\dots (\rho - k + 2)}{k!} \tau^k t^{\rho - k + 1} = \\
%&= x + \sum_{\rho=0}^{p-1} w_\rho \sum_{k=1}^p \frac{\rho!}{k!(\rho - k + 1)!} \tau^k t^{\rho - k + 1}
%\end{align}
Der Induktionsanfang für $j=0$ ist trivial, da per Konstruktion
\begin{align}
x_h(t_0) &= x(t_0).
\end{align}
Angenommen es gilt $x_h(t_{j}) = x(t_{j})$ für ein $j\ge 0$.
Dann folgt
\begin{align}
x_h(t_{j+1}) &= \Psi^{t_{j+1},t_{j}}x_h(t_{j}) = x(t_{j}) + \tau \sum_{l=1}^s b_lk_l,
\end{align}
wobei die Stufen durch
\begin{align}
k_m &= f\kla\left( t_j + c_m\tau \right)
\end{align}
in unserem Fall vollständig entkoppeln.
%Es gilt
%\begin{align}
%\Psi^{t+\tau, t}x = x + \tau \sum_{l=1}^s b_lk_l,
%\end{align}
%wobei die Stufen durch
%\begin{align}
%k_l &= f\kla\left( t + c_l\tau \right)
%\end{align}
%gegeben sind, und somit in unserem Fall vollständig entkoppeln.
Der neue Konsistenzfehler ist dann
\begin{align}
\epsilon_{j+1} &= x_h(t_{j+1}) - x(t_{j+1}) = \underbrace{\epsilon_{j}}_{=\,0} + \Phi^{t_{j+1},t_j} x(t_j) - \Psi^{t_{j+1},t_j} x(t_j).
\end{align}
Es gilt
\begin{align}
\kla\left( \Phi^{t+\tau, t} - \Psi^{t+\tau, t} \right)x &= \sum_{\rho=0}^{p-1} w_\rho \frac{(t+\tau)^{\rho + 1} - t^{\rho + 1}}{\rho + 1} - \tau \sum_{l=1}^s b_l \sum_{\rho = 0}^{p-1} w_\rho \kla\left( t + c_l \tau \right)^\rho = \\
&= \sum_{\rho = 0}^{p-1} w_\rho \klam\left[ \frac{1}{\rho + 1} \sum_{m=0}^{\rho + 1} \binom{\rho+1}{m} t^m \tau^{\rho+1-m} - \frac{t^{\rho + 1}}{\rho + 1} - \tau \sum_{l=1}^s b_l (t + c_l\tau)^\rho \right].
\end{align}
Da die Koeffizienten beliebig sind, muss der Ausdruck in den eckigen Klammern für alle $\rho$ einzeln verschwinden.
Wir können diesen zunächst noch etwas vereinfachen
\begin{align}
[\dots] &= \frac{\tau}{\rho+1} \sum_{m=0}^\rho \frac{(\rho + 1)!}{m!(\rho + 1 - m)!} t^m \tau^{\rho - m} - \tau \sum_{l = 1}^s b_l \sum_{m=0}^\rho \binom{\rho}{m} t^m (c_l\tau)^{\rho - m}  = \\
&= \tau \rho! \sum_{m=0}^\rho \frac{t^m}{m!} \frac{\tau^{\rho - m}}{(\rho - m)!} \klam\left[ \frac{1}{\rho + 1 - m} - \sum_{l=1}^s b_l c_l^{\rho - m} \right].
\end{align}
Da die Schrittweite $\tau$ beliebig ist, und in jedem Term unterschiedliche Potenzen von $\tau$ auftreten, muss wieder jeder Summand verschwinden.
Für alle $m=0,\dots,\rho < p$ muss also gelten
\begin{align}
\frac{1}{m+1} &= \sum_{l=1}^s b_l c_l^m.
\end{align}
Das entspricht gerade der Bedingung an eine Quadraturformel mit $s$ Gewichten bzw. Stützstellen, ein Polynom vom Grad $m$ exakt zu integrieren.
Für $s\ge \frac{p}{2}$ sollte es möglich sein, entsprechende Koeffizienten zu finden.
Damit ist allerdings nur gezeigt, dass prinzipiell ein \textsc{Runge-Kutta}-Verfahren exisitiert, was Polynom exakt ``integriert''.
Das erscheint zum Einen nicht besonders beeindruckend (letztlich beruht das ja wieder auf der ``Magie'' der Quadraturformeln), und zum Anderen nicht wirklich die Aufgabe (wir sollten ja zeigen, dass das für beliebige RK-Verfahren funktioniert).
\\
Die Bedingung wird von den drei Verfahren aus dem Programmier-Teil erfüllt.
Damit haben wir zumindest auch kein offensichtliches Gegenbeispiel zu der Behauptung, dass die Bedingung mit der Quadraturformel tatsächlich notwendig ist.
 

%\begin{thebibliography}{9}
%\bibitem{example} description, \url{https://www.exmapleurl}, day.\,month~2021.
%\end{thebibliography}

\end{document}
